\begin{table}[htbp]
  \centering
  \setlength{\tabcolsep}{4pt}
  \caption{Required sample sizes for the cytisine trial under the design assumptions of the original trial: $\pi_{1} = 0.47$, $\pi_{0} = 0.41$, $\omega_{1} = \omega_{0} = 0.10$, $r = 1$, $\alpha = 0.025$ (one-sided) and target power $= 0.8$. The simple inflation method gives $N_{\mathrm{CC}} = 2{,}388$, which is the sample size the trial planned; it enrolled 2,472.}
  \label{tab:application}
  \begin{threeparttable}
  \begin{tabular}{lccccccrr}
    \toprule
    $\rho$ & $\gamma_{1}$ & $\gamma_{0}$ & $\pi_{1,\mathrm{NRI}}$ & $\pi_{0,\mathrm{NRI}}$ & $\delta_{\mathrm{NRI}}$ & $N_{\mathrm{CC}}$ & $N_{\mathrm{NRI}}$ & RE \\
    \midrule
    $L$ (-0.278) & 0.054 & 0.000 & 0.4646 & 0.4100 & 0.0546 & 2,388 & 2,590 & 0.92 \\
    $0$ & 0.470 & 0.410 & 0.4230 & 0.3690 & 0.0540 & 2,388 & 2,574 & 0.93 \\
    $U$ (0.354) & 1.000 & 0.932 & 0.3700 & 0.3168 & 0.0532 & 2,388 & 2,498 & 0.96 \\
    \bottomrule
  \end{tabular}
  \begin{tablenotes}[flushleft]\footnotesize
    \item $L$ and $U$ denote the lower and upper values of $\rho$ attainable in both groups at once, that is the endpoints of the intersection of the two Prentice intervals; here $L$ is the lower endpoint for the placebo group and $U$ the upper endpoint for the cytisine group. $\gamma_{j} = \Pr(R_{ij} = 1 \mid D_{ij} = 1)$ is the probability that a dropout in group $j$ would have responded, which is the same assumption as $\rho$ expressed on the probability scale; a common $\rho$ implies different $\gamma_{j}$ in the two groups because the response probabilities differ. $\pi_{j,\mathrm{NRI}} = \Pr(R_{ij} = 1, D_{ij} = 0)$ is the response probability an NRI analysis observes, and $\delta_{\mathrm{NRI}} = \pi_{1,\mathrm{NRI}} - \pi_{0,\mathrm{NRI}}$ is the effect it is powered against; the corresponding full-data effect is $\delta_{\mathrm{Full}} = 0.06$. RE $= N_{\mathrm{CC}} / N_{\mathrm{NRI}}$.
  \end{tablenotes}
  \end{threeparttable}
\end{table}
